%% kapa: exam=1 round=20 year=2009 period=2 form=A subject=accounting subject_ko=회계학 seq=31 number=31
%% source: Q-Net 감정평가사 기출문제 (cst00309) artlSeq=5024178 file=extracted/1008/5024178/2009년 제20회 감정평가사 1차시험 기출문제/2교시.hwp
%% conversion: kordoc 4.24.51 -> tools/gen1.py (auto); review=proofread-sonnet; flags=none
\begin{question}{31}
(주)한국은 실제원가계산을 적용하고 있으며, 20×9년 1월의 월초 및 월말 재고자산 금액은 다음과 같다.

{\small\centering
\begin{tabular}{@{}l r r r@{}}
\hline
 & \multicolumn{1}{c}{직접재료} & \multicolumn{1}{c}{재공품} & \multicolumn{1}{c}{제품} \\
\cline{2-2}\cline{3-3}\cline{4-4}
1월초 & ￦25,000 & ￦30,000 & ￦40,000 \\
1월말 & 15,000 & 20,000 & 25,000 \\
\hline
\end{tabular}\par}
1월중에 가공원가는 ￦230,000이 발생하였으며, 재공품계정의 차변합계 금액은 ￦330,000이었다. 20×9년 1월의 직접재료 구입액과 매출원가는 각각 얼마인가? \par

{\small\centering
\begin{tabular}{@{}ccc@{}}
 & 직접재료 구입액 & 매출원가 \\
\cline{2-2}\cline{3-3}
① & ￦70,000 & ￦325,000 \\
② & ￦70,000 & ￦350,000 \\
③ & ￦65,000 & ￦325,000 \\
④ & ￦60,000 & ￦325,000 \\
⑤ & ￦60,000 & ￦350,000 \\
\end{tabular}\par}
\end{question}
\answer{4}
